\chapter{Trabalhos Relacionados}\label{cap:relacionados}

O processamento de linguagem natural clínico consolidou-se como área de pesquisa, com aplicações em extração de fenótipos, codificação diagnóstica, farmacovigilância e apoio à decisão. A representação do conhecimento clínico em grafos, ligando pacientes, diagnósticos, exames e terapias, tem sido explorada para viabilizar consultas estruturadas e raciocínio sobre trajetórias de doença. A combinação de LLMs com grafos, na forma de GraphRAG, é desenvolvimento recente nesse cenário, e os trabalhos disponíveis diferem quanto ao tipo de dado, à arquitetura e à forma de avaliação. Este capítulo revisa esses trabalhos, a literatura sobre engenharia de \textit{prompt} em processamento de linguagem natural clínico e situa a lacuna que motiva a presente pesquisa.

\section{GraphRAG em Registros Clínicos e Saúde}\label{sec:relacionados-graphrag-saude}

\citet{baptista2025} aplicaram GraphRAG a prontuários eletrônicos representados no padrão FHIR, com uma arquitetura própria de expansão de vizinhança a um salto (\textit{1-hop}) sobre o grafo, e avaliaram cinco LLMs distintos no papel de gerador. \citet{baptista2026} deu continuidade a esse trabalho. Os dados utilizados são sintéticos, gerados pelo simulador Synthea, e a avaliação apoia-se em métricas de abrangência e diversidade das respostas, atribuídas por um LLM no papel de juiz. A presente pesquisa difere por operar sobre prontuários reais em texto livre, e não sobre registros estruturados sintéticos, e por adotar avaliação com \textit{ground truth} anotado por especialista e métricas objetivas de classificação.

A geração de respostas ancoradas em evidência extraída do prontuário passou, recentemente, a contar com avaliação padronizada. A tarefa compartilhada ArchEHR-QA 2025 \citep{soni2025} pediu que os sistemas respondessem a perguntas de pacientes com respostas fundamentadas nas evidências clínicas centrais de seus próprios prontuários eletrônicos, e reuniu 29 equipes e 75 sistemas. As arquiteturas submetidas incluíram abordagens que selecionam as evidências mais relevantes antes de gerar a resposta, com LLMs proprietários e de pesos abertos, adaptados por \textit{fine-tuning} ou por aprendizado com poucos exemplos (\textit{few-shot}). A estratégia de selecionar evidência antes de gerar é o mesmo princípio que fundamenta a geração aumentada por recuperação, discutida no capítulo seguinte; a tarefa, porém, é de geração de resposta livre, ao passo que este trabalho trata de classificação multirrótulo de diagnósticos registrados no prontuário.

O sistema KARE \citep{kare2025} emprega recuperação baseada em comunidades para tarefas de predição clínica sobre prontuários. Para obter comunidades em diferentes granularidades, os autores executam o algoritmo de Leiden em múltiplos níveis, contornando a granularidade fixa adotada pela implementação do GraphRAG à época. A tarefa de KARE é de predição de desfechos, ao passo que este trabalho trata da identificação de um diagnóstico já registrado no prontuário.

Outra vertente aplica o GraphRAG não ao prontuário individual, mas a bases de conhecimento médico. O MedGraphRAG \citep{wu2025} foi concebido para que LLMs gerem respostas médicas baseadas em evidência, com o objetivo declarado de aumentar a segurança e a confiabilidade das respostas produzidas a partir de dados médicos privados. Para isso, propõe uma construção de grafo em triplas, que vincula os documentos do usuário a fontes médicas de credibilidade estabelecida, e uma estratégia de recuperação denominada U-Retrieval, que combina uma busca descendente (\textit{top-down}), voltada à indexação precisa, a um refinamento ascendente (\textit{bottom-up}) da resposta, voltado ao equilíbrio do contexto. Validado em nove \textit{benchmarks} de perguntas e respostas médicas e em dois conjuntos de verificação de fatos em saúde, o sistema superou modelos do estado da arte. A combinação entre uma recuperação ampla, que percorre o conhecimento de cima para baixo, e um refinamento localizado expressa a mesma tensão entre abrangência e especificidade que este trabalho examina empiricamente, ao comparar modos de busca e níveis de comunidade. \citet{wang2026}, por sua vez, investigam o GraphRAG em medicina interdisciplinar, na qual a evidência provém de sistemas de conhecimento incomensuráveis, como a medicina tradicional chinesa e a medicina ocidental. Os autores argumentam que a recuperação por similaridade em um espaço latente compartilhado tende a unificar em excesso conceitos heterogêneos, o que provoca um colapso semântico, com alinhamentos espúrios, mistura de procedências e cadeias de raciocínio em múltiplos saltos rompidas; em vez de unificar as ontologias, constroem um grafo de ontologia dupla conectado por pontes tipadas, com ganhos de até 22,2\% de acurácia sobre o \textit{baseline} OG-RAG. O problema que descrevem é o simétrico de um dos observados neste trabalho: enquanto ali conceitos distintos são indevidamente fundidos, aqui formas sinônimas de uma mesma doença permanecem dispersas em entidades separadas, como se discute no capítulo de resultados. Ambos os trabalhos, de todo modo, avaliam o GraphRAG em perguntas e respostas sobre conhecimento médico, e não na extração de diagnósticos a partir de texto clínico real.

\citet{rakuljic2026} compararam RAG tradicional e GraphRAG sobre dados corporativos, em tarefas de perguntas e respostas, e relataram superioridade do RAG tradicional na maioria das tarefas avaliadas. O domínio e o tipo de documento diferem dos deste trabalho, mas o resultado dialoga com a discussão sobre contextos longos e degradação da atenção retomada no capítulo seguinte.

\citet{esteves2025} avaliou GraphRAG no domínio jurídico em português europeu e documentou dificuldades associadas ao processamento de línguas sub-representadas, como a menor qualidade de modelos de extração e de \textit{embeddings}. Esta pesquisa compartilha o desafio linguístico, no par formado pelo português brasileiro e pelo domínio clínico, com o agravante da linguagem telegráfica dos prontuários.

\section{Engenharia de Prompt em PLN Clínico}\label{sec:relacionados-prompt}

Além dos trabalhos que aplicam GraphRAG especificamente, a literatura de processamento de linguagem natural clínico oferece evidência direta sobre o papel do \textit{prompt} na extração de informação a partir de texto livre. \citet{jesus2026} avaliaram cinco LLMs (GPT-4o mini, DeepSeek-V3, Mixtral-8x7B, LLaMA 8B e Qwen-32B) na extração de dados clínicos de prontuários ambulatoriais em português brasileiro, nos domínios de cardiologia e infectologia, e relataram acurácia balanceada de até 0,91, com desempenho mais baixo em condições pouco documentadas ou ambíguas, como neoplasias, frente a condições bem registradas, como diabetes. \citet{pereira2026} compararam métodos determinísticos, modelos generativos com engenharia de \textit{prompt}, modelos ajustados por \textit{fine-tuning} e modelos de raciocínio sobre notas clínicas em português, em tarefas de complexidade semântica crescente, e constataram que variáveis mais simples são bem tratadas por métodos lexicais determinísticos, ao passo que variáveis ambíguas exigem modelos mais sofisticados, a um custo computacional maior. Para o português clínico especificamente, esses dois trabalhos reforçam o achado mais geral de que o desempenho de um LLM em extração clínica depende da formulação da tarefa no \textit{prompt}, e não apenas da capacidade do modelo em si.

O mesmo padrão aparece em outras línguas sub-representadas na pesquisa em PLN clínico. \citet{lilli2025} compararam LLMs generativos e modelos baseados em BERT na classificação de laudos clínicos reais em italiano, distribuídos em oito categorias diagnósticas da área de reumatologia, sob três estratégias de adaptação: \textit{prompt} sem exemplos (\textit{zero-shot}), aprendizado em contexto (\textit{in-context learning}) e ajuste fino. O aprendizado em contexto superou o \textit{zero-shot}, sobretudo nos modelos Mixtral e Gemma; o ajuste fino de BERT obteve o melhor desempenho global, mas foi superado pelo aprendizado em contexto em diagnósticos específicos, como os renais e os sistêmicos, o que levou os autores a apontar o \textit{prompting} como alternativa viável quando dados clínicos rotulados são escassos. O paralelo com este trabalho é direto: também aqui se classifica texto clínico real, em língua diferente do inglês, em um conjunto fixo de oito categorias diagnósticas, e por meio de \textit{prompting}, sem ajuste fino de um modelo dedicado.

No caso específico do GraphRAG, \citet{ye2026} relatam que a aplicação do pipeline padrão a prontuários eletrônicos produz grafos com granularidade excessiva, o que degrada a qualidade da recuperação, e corrigem o problema por meio de engenharia de \textit{prompt} adaptada ao domínio médico, além de propor a verificação da resposta contra conhecimento estruturado como defesa contra alucinação. O achado dialoga diretamente com três aspectos deste trabalho: o ajuste dos \textit{prompts} de indexação ao domínio, descrito na Seção~\ref{sec:met-grafo}; o refinamento do \textit{prompt} de extração, descrito na Seção~\ref{sec:met-prompt} e motivado por problemas concretos identificados ao longo das rodadas experimentais; e a investigação da alucinação do modo \textit{drift} relatada no capítulo de resultados.

\section{Lacuna de Pesquisa}\label{sec:relacionados-lacuna}

A revisão desses trabalhos evidencia uma lacuna. Não há, na literatura consultada, avaliação do GraphRAG sobre prontuários clínicos reais em português brasileiro, em domínio médico especializado, com \textit{ground truth} anotado por humano e comparação sistemática entre os modos de busca e os níveis de comunidade. É essa lacuna que a presente pesquisa se propõe a preencher.

Antes de detalhar a metodologia que preenche essa lacuna, o próximo capítulo reúne os fundamentos teóricos necessários à sua compreensão.
